\chapter{Core Combinatory Calculus}

In this chapter, we formalize the syntax and reduction semantics of the ISAR combinatory logic kernel. The core ISAR calculus is a substrate-independent combinator system that supports S-K combinator reductions, parallel reductions, and normalization.

\section{Syntax and Reduction Semantics}

The syntax consists of standard S-K combinators and the specialized ISAR terms.

\begin{definition}[S-K Term Syntax]
\lean{ISAR.SK}
The inductive type of basic S-K combinator expressions is defined by:
\begin{align*}
t, u \in \text{SK} &::= \mathbf{S} \mid \mathbf{K} \mid t\ u
\end{align*}
\end{definition}

\begin{definition}[ISAR Term Syntax]
\lean{ISAR.ITerm}
The core term syntax of the ISAR kernel extends basic combinators with norms, constants, and custom operations:
\begin{align*}
t, u \in \text{ITerm} &::= \mathbf{norm} \mid \mathbf{s_s} \mid \mathbf{konst} \mid \mathbf{app}\ t\ u
\end{align*}
\end{definition}

We define the reduction rules for both systems.

\begin{definition}[SKStep Reduction]
\lean{ISAR.SKStep}
The reduction relation on SK terms represents the standard combinatory logic reduction rules:
\begin{itemize}
    \item $\mathbf{K}\ x\ y \to x$
    \item $\mathbf{S}\ x\ y\ z \to (x\ z)(y\ z)$
\end{itemize}
\end{definition}

\begin{definition}[IStep Reduction]
\lean{ISAR.IStep}
The single-step reduction relation on ISAR terms defines the operational semantics of norms and constants.
\end{definition}

\begin{theorem}[Confluence of IRed]
\lean{ISAR.IRed_confluence}
\uses{ISAR.ITerm, ISAR.IStep}
The reduction relation $\text{IRed}$ (the reflexive transitive closure of $\text{IStep}$) satisfies the Church-Rosser confluence property:
$$\forall t, u_1, u_2, \quad t \to^* u_1 \land t \to^* u_2 \implies \exists v, \quad u_1 \to^* v \land u_2 \to^* v$$
\end{theorem}

\begin{theorem}[Unique Normal Forms]
\lean{ISAR.isar_fragment_unique_normal_forms}
\uses{ISAR.IRed_confluence}
If a term $t$ reduces to two normal forms $n_1$ and $n_2$, then $n_1 = n_2$.
\end{theorem}

\section{The Invariant Layer}

Using the uniqueness of normal forms, we define the quotient space representing behavioral equivalence classes of terms.

\begin{definition}[ISKSubtype]
\lean{ISAR.ISKSubtype}
\uses{ISAR.ITerm}
The subtype of ISAR terms belonging to the S-K fragment:
$$\text{ISKSubtype} := \{ t : \text{ITerm} \mid \text{ISKTerm } t \}$$
\end{definition}

\begin{definition}[Operational Equivalence]
\lean{ISAR.OperEq}
\uses{ISAR.ISKSubtype}
Two terms are operationally equivalent if they behave identically under reduction:
$$t \sim_{op} u \iff \text{OperEq } t\ u$$
\end{definition}

\begin{definition}[Invariant Layer]
\lean{ISAR.InvariantLayer}
\uses{ISAR.OperEq}
The quotient type of \texttt{ISKSubtype} modulo \texttt{OperEq}:
$$\text{InvariantLayer} := \text{Quotient } (\text{operEqSetoid})$$
\end{definition}

\begin{definition}[Invariant Layer Application]
\lean{ISAR.InvariantLayer.app}
\uses{ISAR.InvariantLayer}
The application operator lifted to the quotient space:
$$\text{app} : \text{InvariantLayer} \to \text{InvariantLayer} \to \text{InvariantLayer}$$
\end{definition}

\begin{definition}[Canonical Representative]
\lean{ISAR.InvariantLayer.canonical_rep}
\uses{ISAR.InvariantLayer}
A noncomputable function mapping each equivalence class to its canonical representative in \texttt{ISKSubtype}.
\end{definition}
